\chapter{Tecnologie}
In questo capitolo vengono descritte le tecnologie adottate per lo sviluppo dell'applicazione Easy Notes, insieme alle motivazioni che hanno portato a tali scelte.
\section{Stack}
L'applicazione è sviluppata utilizzando lo stack \b{MEVN}, per cui:
\begin{itemize}
    \item \textbf{MongoDB} come database NoSQL per la gestione dei dati;
    \item \textbf{Express.js} come framework backend per Node.js;
    \item \textbf{Vue.js} come framework frontend per la costruzione dell'interfaccia utente;
    \item \textbf{Node.js} come runtime environment per l'esecuzione del codice JavaScript lato server.
\end{itemize} 
La scelta viene motivata direttamente dalle richieste del progetto d'esame.

Lo stack MEVN viene messo in pratica attraverso l'uso di Microsoft Azure come piattaforma di cloud computing; infatti l'applicazione viene distribuita su un'istanza di \textbf{Azure App Service} che ospita sia il backend che il frontend, mentre il database MongoDB è gestito tramite \textbf{Azure Cosmos DB}.

\section{Librerie e Framework}
Per lo sviluppo dell'applicazione sono state utilizzate le seguenti librerie e framework principali:
\begin{itemize}
    \item \textbf{GridFS}: per la gestione e l'archiviazione di file di grandi dimensioni all'interno di MongoDB, nel nostro caso cruciale per poter salvare i file pdf all'interno del database senza dover ricorrere a soluzioni esterne tramite URL;
    \item \textbf{Mongoose}: per la modellazione degli oggetti MongoDB in ambiente Node.js, semplificando l'interazione con il database attraverso uno schema definito;
    \item \textbf{Axios}: per effettuare richieste HTTP dal frontend Vue.js al backend Express.js, facilitando la comunicazione tra client e server;
    \item \textbf{Vue Router}: per la gestione della navigazione tra le diverse pagine dell'applicazione Vue.js, permettendo una struttura a singola pagina (SPA);
    \item \textbf{Cerebras SDK}: per l'integrazione con i modelli di intelligenza artificiale offerti da Cerebras, utilizzati per riassumere note e pdf tramite modello llama-3.3-70b
\end{itemize}

\section{Strumenti di Sviluppo}
Per lo sviluppo dell'applicazione sono stati utilizzati i seguenti strumenti:
\begin{itemize}
    \item \textbf{Visual Studio Code}: come ambiente di sviluppo integrato (IDE) per la scrittura e la gestione del codice sorgente;
    \item \textbf{Postman}: per testare e debug le API RESTful sviluppate con Express.js, facilitando la verifica delle funzionalità del backend;
    \item \textbf{Git}: per il controllo di versione del codice sorgente, permettendo una gestione efficiente delle modifiche e la collaborazione tra sviluppatori;
    \item \textbf{Docker}: per la containerizzazione dell'applicazione, semplificando il processo di deployment su Azure App Service.
\end{itemize}
